\documentclass[../../../include/open-logic-section]{subfiles}

\begin{document}

\olfileid{sth}{choice}{countablechoice}
\olsection{তালিকায়নযোগ্য চয়ন}

চয়ন বা সুক্রমবিন্যাস কোনোটি না ধরেই সহজে প্রমাণ
করা যায় যে:

\begin{lem}[$\Zminus$-এ]
প্রত্যেক সসীম সেটের একটি চয়ন-অপেক্ষক আছে।
\end{lem}

\begin{proof}
$a = \{b_1, \ldots, b_n\}$ ধরি। সরলতার জন্য
প্রত্যেক $b_i\neq\emptyset$ ধরি। তাহলে এমন
$c_1,\ldots,c_n$ আছে যে
$c_1\in b_1,\ldots,c_n\in b_n$।
\olref[z][pairs]{prop:pairsconsequences} অনুযায়ী
$\{\langle b_1,c_1\rangle,\ldots,
\langle b_n,c_n\rangle\}$ সেটটি আছে।
এটি $a$-র একটি চয়ন-অপেক্ষক।
\end{proof}

কিন্তু অসীম সেটের কথায় গেলেই পরিস্থিতি অস্পষ্ট।
যেমন, আগের ফলটির এই ``ন্যূনতম'' বিস্তার ধরো:

\begin{defish}
\emph{তালিকায়নযোগ্য চয়ন।} প্রত্যেক
\emph{তালিকায়নযোগ্য} সেটের একটি চয়ন-অপেক্ষক আছে।
\end{defish}

এটি চয়নের একটি বিশেষ ক্ষেত্র। দেখা যায়,
নীতিটি বহুবার প্রয়োগ করা হয়েছে, অথচ প্রয়োগকারীরা
যে তা বুঝেছিলেন, এমন স্পষ্ট লক্ষণ নেই। দুটি
উদাহরণ দেখি।\footnote{\citet[\S9.4]{Potter2004}
ও লুকা ইনকুরভাতির সূত্রে।}

\begin{ex}
সহজে মনে হতে পারে: যেকোনো সেট $A$-র জন্য
হয় $\cardle{\omega}{A}$, নয় কোনো
$n\in\omega$-র জন্য $\cardeq{A}{n}$।
প্রত্যেক সেট হয় সসীম, নয় অসীম—এই স্বজ্ঞার
একটি রূপ এটি। কান্টরসহ বহু গণিতবিদ দাবিটি
প্রমাণ ছাড়াই করেছিলেন। আমরা সতর্ক বলে
\olref[cardinals][classing]{generalinfinitycharacter}-এ
এটি প্রমাণ করেছি। কিন্তু সেই প্রমাণে আমরা $\ZFC$-এ
কাজ করছিলাম: যেকোনো $A$ সুক্রমবিন্যাসযোগ্য
এবং তাই $\card{A}$ আছে, ধরে নিয়েছিলাম।
অর্থাৎ স্পষ্টভাবেই চয়ন ধরেছিলাম।

আসলে \citet{Dedekind1888} দাবিটির নিজের
প্রমাণ দিয়েছিলেন:

\begin{thm}[$\Zminus + \text{তালিকায়নযোগ্য চয়ন}$-এ]
যেকোনো $A$-র জন্য হয়
$\cardle{\omega}{A}$, নয় কোনো $n\in\omega$-র
জন্য $\cardeq{A}{n}$।
\end{thm}

\begin{proof}
প্রত্যেক $n\in\omega$-র জন্য
$\cardneq{A}{n}$ ধরি। তাহলে বিশেষত প্রত্যেক
$n<\omega$-র জন্য $A$-র এমন একটি উপসেট
$A_n\subseteq A$ আছে যার সদস্যসংখ্যা ঠিক
$\cardexpo{2}{n}$। এই $A_0,A_1,A_2,\ldots$
অনুক্রম দিয়ে প্রত্যেক $n$-এর জন্য নিই:
\[
	B_n = A_n \setminus \bigcup_{i < n} A_i.
\]
এবার লক্ষ করো:
% BN-SRC-470: the union ranges over A_i, not a repeated A_n.
\begin{align*}
	\card{\bigcup_{i < n}A_i}
	&\leq \card{A_0} + \card{A_1} + \ldots + \card{A_{n-1}}\\
	&=1 + 2 + \ldots + 2^{n-1}\\
	& = 2^n - 1\\
	& < 2^n = \card{A_n}
\end{align*}
তাই প্রত্যেক $B_n$-এর অন্তত একটি সদস্য
$c_n$ আছে। তদুপরি $B_n$-গুলি জোড়ায় জোড়ায়
বিচ্ছিন্ন। ফলে $c_n=c_m$ হলে $n=m$।
আবার প্রত্যেক $c_n\in A$। তাই $f(n)=c_n$
একটি $\omega\to A$ অন্তঃস্থাপন।
% BN-SRC-823 TARGET-CORRECTION-BEGIN
\par\noindent\textbf{পাঠ-সংশোধন।} এখানে $A$-কে আগে সুক্রমবিন্যস্ত করা হয়নি। $A$ কোনও
সসীম সেটের সমসংখ্যক নয় বলে প্রত্যেক সসীম সংখ্যক পৃথক
সদস্য নেওয়া যায়: সসীম অংশ বাদ দিলে আর সদস্য না থাকলে
$A$-ই সসীম হতো। ধরি $S_n$ হলো $A$-র ঠিক $2^n$ সদস্যের
উপসেটগুলির অশূন্য সেট। এগুলি $\Pow{A}$-র উপসেট, তাই
তাদের অনুক্রমের graph এবং সদস্য-সেটকে বদ্ধ ঘাতসেটে
পৃথকীকরণে গড়া যায়। প্রথম তালিকায়নযোগ্য চয়নে
$A_n\in S_n$, দ্বিতীয়টিতে $c_n\in B_n$।
শেষ graph $\Setabs{\tuple{n,c_n}\in\omega\times A}{n\in\omega}$
বদ্ধ পৃথকীকরণেই সেট; প্রতিস্থাপন গোপনে লাগে না।
সসীম সেটের আগের লেমায় পৃথক সদস্যগুলি একবার করে তালিকাভুক্ত
করে শূন্য সেট বাদ দিতে হয়; সসীম বহু অস্তিত্বের সাক্ষী একসঙ্গে
নেওয়া যুক্তিবিদ্যার সসীম অনুমিতি, চয়নের স্বতঃসিদ্ধ নয়।
এই উপপাদ্য সাধারণ ``সসীম নয়''-এর চেয়ে বেশি কিছু বলে:
তালিকায়নযোগ্য অসীম উপসেট তথা ডেডেকিন্ড-অসীমতা দেয়।
% BN-SRC-823 TARGET-CORRECTION-END
\end{proof}
\noindent
দেদেকিন্ড বলেননি যে তিনি তালিকায়নযোগ্য চয়ন
ব্যবহার করেছেন। কিন্তু \emph{তুমি} কি ব্যবহারটি
ধরতে পেরেছিলে? আবার দেখো। (সত্যিই:
\emph{আবার দেখো}।)

প্রমাণে তালিকায়নযোগ্য চয়ন দুবার লেগেছে।
প্রথমবার $A_0$, $A_1$, $A_2$, \dots\@ সেটগুলির
অনুক্রম পেতে। দ্বিতীয়বার প্রত্যেক $B_n$ থেকে
$c_n$ বেছে নিতে। আর চয়নের এই ব্যবহার
বাদ দেওয়া যায় না। \citet[p.~138]{Cohen1966}
দেখিয়েছিলেন, চয়নের কোনো রূপ না ধরলে ফলটি
ব্যর্থ হতে পারে। অর্থাৎ $\ZF$-র সঙ্গে এমন
সেটের অস্তিত্ব সংগতি-সম্মত, যার আকার
$\omega$-র সঙ্গে \emph{অতুলনীয়}।
\end{ex}

\begin{ex}
\citeyear{Cantor1878}-এ কান্টর বলেছিলেন,
তালিকায়নযোগ্য সংখ্যক তালিকায়নযোগ্য সেটের
সংযোগও তালিকায়নযোগ্য। প্রমাণ দেননি;
হয়তো তাঁর কাছে ফলটি স্পষ্ট ছিল। আমরা সতর্ক
বলে \olref[card-arithmetic][simp]{kappaunionkappasize}-এ
আরও সাধারণ ফল প্রমাণ করেছি। কিন্তু সেখানে
স্পষ্টভাবেই চয়ন ধরেছি। কম সাধারণ ফলটির
প্রমাণেও তালিকায়নযোগ্য চয়ন লাগে।

\begin{thm}[$\Zminus + \text{তালিকায়নযোগ্য চয়ন}$-এ]
প্রত্যেক $n\in\omega$-র জন্য $A_n$
তালিকায়নযোগ্য হলে
$\bigcup_{n<\omega}A_n$-ও তালিকায়নযোগ্য।
\end{thm}

\begin{proof}
সাধারণত্ব না হারিয়ে প্রত্যেক
$A_n\neq\emptyset$ ধরি। তাই প্রত্যেক
$n\in\omega$-র জন্য একটি !!a{surjection}
$f_n\colon\omega\to A_n$ আছে।
$f\colon\omega\times\omega\to
\bigcup_{n<\omega}A_n$-কে
$f(m,n)=f_n(m)$ দিয়ে সংজ্ঞায়িত করি।
\olref[sfr][siz][zigzag]{natsquaredenumerable}
অনুযায়ী $\omega\times\omega$
তালিকায়নযোগ্য, আর $f$ একটি
!!a{surjection}; তাই ফলটি মেলে।
% BN-SRC-824 TARGET-CORRECTION-BEGIN
\par\noindent\textbf{পাঠ-সংশোধন।} শূন্য সদস্য-সেট বাদ দেওয়ার ধাপটি এভাবে করা যায়। সংযোগ
$U=\bigcup_{n<\omega}A_n$ শূন্য হলে দাবিটি সরাসরি সত্য।
নইলে একটি $u\in U$ স্থির করে প্রতিটি শূন্য $A_n$-কে
$\{u\}$ দিয়ে বদলাই; সংযোগ বদলায় না, সব সদস্য-সেট অশূন্য
ও তালিকায়নযোগ্য হয়। প্রত্যেকটির উপর $\omega$ থেকে
সমাপতিত চিত্রণগুলির সেট $\Pow{\omega\times U}$-তে বদ্ধ;
এদের তালিকায়নযোগ্য পরিবার থেকে চিত্রণগুলি বেছে নিই।
সব নির্বাচিত graph-কে বদ্ধ সেটে পৃথকীকরণে গড়া যায়,
তাই এই $\Zminus$-প্রমাণেও প্রতিস্থাপন অনুক্ত নয়।
$\omega\times\omega$-র স্থির তালিকা সংযোজন করে $U$-র
একটি তালিকা পাই। উল্টো একৈক চিত্রণ চাইলে প্রত্যেক সদস্যের
ক্ষুদ্রতম তালিকা-সূচক নিই; এর জন্য নতুন চয়ন লাগে না।
% BN-SRC-824 TARGET-CORRECTION-END
\end{proof}
\noindent
তালিকায়নযোগ্য চয়ন কোথায় লেগেছে,
ধরতে পেরেছিলে? $f_0$, $f_1$, $f_2$, \dots
অপেক্ষকগুলির অনুক্রম বেছে নিতে লেগেছে।\footnote{
\olref[card-arithmetic][simp]{kappaunionkappasize}-এও
একইভাবে চয়ন লেগেছিল; সেখানে বলেছিলাম,
``প্রত্যেক $\beta\in\cardfont{a}$-র জন্য একটি
!!a{injection} $f_\beta$ স্থির করি।''}
আবার, চয়নের কোনো নীতি না থাকলে এই ফলও
ব্যর্থ হতে পারে। বিশেষত
\citet{FefermanLevy1963} দেখিয়েছিলেন,
$\ZF$-র সঙ্গে এ কথাটি সংগতি-সম্মত যে
তালিকায়নযোগ্য সেটগুলির তালিকায়নযোগ্য
সংযোগটি $\Real$-এর সমসংখ্যক।
% BN-SRC-825 TARGET-CORRECTION-BEGIN
\par\noindent\textbf{পাঠ-সংশোধন।} এই ফলটি $\ZF$-র সঙ্গতি-সাপেক্ষ একটি মডেল সম্বন্ধে:
সেখানে বাস্তব সংখ্যার সেট নিজেই তালিকায়নযোগ্য সেটগুলির
তালিকায়নযোগ্য সংযোগ, অথচ তালিকায়নযোগ্য নয়। চয়ন ছাড়া
আগের প্রারম্ভিক ক্রমসংখ্যা-রীতিতে $\beth_1$-কে বাস্তবগুলির
অঙ্কবাচক সংখ্যা ধরে নেওয়া যায় না। তাই এখানে সমসংখ্যকতা
দিয়ে কথা বলা হয়েছে; $\Real\approx\Pow{\omega}$-র
চয়নবিহীন ফলটি বহাল থাকে। যদি এই মডেলে $\Real$ সুক্রমবিন্যাসযোগ্য
হতো, তবে $\Real^\omega\approx\Real$ থেকে তার তালিকায়নযোগ্য
তালিকাগুলিরও সুক্রম পাওয়া যেত। প্রতিটি সদস্য-সেটের ক্ষুদ্রতম
তালিকা বেছে সংযোগ তালিকায়নযোগ্য হতো—বিরোধ।
% BN-SRC-825 TARGET-CORRECTION-END
বিষয়টি রাসেলের এই মজার উদাহরণে আরও
স্পষ্ট হয়:
\begin{quote}
  এক ধনকুবের যতবার এক জোড়া বুট কিনতেন,
  ততবার এক জোড়া মোজাও কিনতেন; অন্য সময়
  কখনো মোজা কিনতেন না। দুটোই এত কিনতেন যে
  শেষ পর্যন্ত তাঁর $\aleph_0$ জোড়া বুট এবং
  $\aleph_0$ জোড়া মোজা হলো\dots\@
  বুটের ডান ও বাঁ আলাদা করা যায়; তাই প্রত্যেক
  জোড়া থেকে একটি বেছে নিতে পারি—সব ডান বুট
  বা সব বাঁ বুট। কিন্তু মোজার ক্ষেত্রে এমন
  কোনো নির্বাচননীতি নিজে থেকে পাওয়া যায় না।
  গুণের স্বতঃসিদ্ধ [অর্থাৎ কার্যত চয়ন] না ধরলে,
  প্রত্যেক জোড়া থেকে একটি করে মোজা নিয়ে
  কোনো শ্রেণি আছে, তা নিশ্চিত করে বলতে পারি না।
  \citep[p.~126]{Russell1919}
\end{quote}
সংক্ষেপে, তালিকায়নযোগ্য সংখ্যক মোজার জোড়া
থাকলে মোজা (কেবল) তালিকায়নযোগ্য সংখ্যক—এটি
প্রমাণে চয়নের কোনো রূপ লাগে। বস্তুত
উপযুক্ত চয়ননীতি না ধরলে তালিকায়নযোগ্য সেটগুলির তালিকায়নযোগ্য
সংযোগ তালিকায়নযোগ্য নাও হতে পারে।
% BN-SRC-826 TARGET-CORRECTION-BEGIN
\par\noindent\textbf{পাঠ-সংশোধন।} তালিকায়নযোগ্য চয়ন এই সংযোগ-ফলের জন্য যথেষ্ট। কিন্তু
দেওয়া proof-এ নীতিটি ব্যবহৃত হয়েছে বলে গোটা তালিকায়নযোগ্য
চয়ন এই ফলের সমতুল্য—এমন উল্টো implication প্রমাণ হয়নি।
মোজার জোড়াগুলির ক্ষেত্রেও আগে সংযোগের একটি তালিকা থাকলে
প্রতিটি জোড়ার ক্ষুদ্রতম তালিকা-সূচকের মোজা বেছে নেওয়া যায়;
এই বিশেষ পরিবারে আলাদা চয়ন আর লাগে না। এখানে যথেষ্ট
নীতি, বিশেষ পরিবারে ব্যবহৃত চয়ন, এবং সেই ফলকে $\ZF$-তে
প্রমাণ করা না যাওয়া—তিনটি দাবির পরিধি আলাদা।
% BN-SRC-826 TARGET-CORRECTION-END
\end{ex}

শিক্ষাটি হলো, প্রয়োগকারীরা অনেক সময় না বুঝেই
তালিকায়নযোগ্য চয়ন ব্যবহার করেছেন। দার্শনিক
প্রশ্ন: তালিকায়নযোগ্য চয়নের \emph{ন্যায্যতা}
কীভাবে দেখাব?

স্বজ্ঞাভিত্তিক একটি চেষ্টা অসীমসংখ্যক ধাপের
একটি কাজের কথা বলতে পারে। ধরো, প্রথম পছন্দটি
$\nicefrac{1}{2}$ মিনিটে, দ্বিতীয়টি
$\nicefrac{1}{4}$ মিনিটে, \dots, $n$-তমটি
$\nicefrac{1}{2^n}$ মিনিটে, \dots\@ করি।
তাহলে $1$ মিনিটের মধ্যেই $\omega$-দৈর্ঘ্যের
পছন্দের অনুক্রম এবং একটি চয়ন-অপেক্ষক তৈরি
হয়ে যাবে।

কিন্তু এমন ভাবনাপরীক্ষা সত্যিই কী দেখায়?
প্রথমত, এটি ``বেছে নেওয়া'' কথাটিকে বেশ
আক্ষরিক অর্থে ধরে। দ্বিতীয়ত, এটি গণিতকে
অধিবিদ্যাগত সম্ভাবনার সঙ্গে বেঁধে ফেলে।

আরও গুরুত্বপূর্ণ, এতে \emph{সরাসরি চয়ন}-এর
ন্যায্যতা মেলে না; বড়জোর \emph{তালিকায়নযোগ্য
চয়ন}-এর ন্যায্যতা মিলতে পারে। \emph{প্রত্যেক}
সেটের চয়ন-অপেক্ষক চাইলে ``যেকোনো
ক্রমসংখ্যা-দৈর্ঘ্যের অসীমধাপী কাজ'' সম্পন্ন করার
সামর্থ্য লাগবে। সোজা কথায়, ধারণাটি হাস্যকর।

\end{document}
